\documentclass[10pt,a4paper]{article}
\usepackage[T1]{fontenc}
\usepackage[utf8]{inputenc}
\usepackage[spanish,es-tabla,es-nodecimaldot]{babel}
\usepackage{mathpazo}
\usepackage[scaled=0.9]{beramono}
\usepackage{microtype}
\usepackage[a4paper,margin=2.4cm]{geometry}
\usepackage{amsmath,amssymb}
\usepackage{booktabs}
\usepackage{tabularx}
\usepackage{caption}
\usepackage{fancyhdr}
\usepackage{titlesec}
\usepackage{tikz}
\usetikzlibrary{arrows.meta,positioning,babel}
\usepackage[hidelinks]{hyperref}

\captionsetup{font=small,labelfont=bf}
\titlespacing*{\section}{0pt}{2.2ex plus .5ex}{1.0ex}
\titlespacing*{\subsection}{0pt}{1.8ex plus .4ex}{0.8ex}
\setlength{\parskip}{0.25ex}
\pagestyle{fancy}
\fancyhf{}
\lhead{\small OII464: segundo informe}
\rhead{\small José Lara Arce}
\cfoot{\thepage}

\DeclareMathOperator*{\maxop}{m\acute{a}x}
\DeclareMathOperator*{\minop}{m\acute{\imath}n}
\newcommand{\Cmax}{C_{\text{máx}}}
\newcommand{\dc}{{,}}% coma decimal en modo matemático

\begin{document}

\begin{center}
{\LARGE\bfseries Planificación quirúrgica con dPSO\\[2pt]
y control aprendido de la perturbación en dos niveles\par}
\vspace{8pt}
{\large José Lara Arce}\\[2pt]
{\small Escuela de Ingeniería Informática, Pontificia Universidad Católica de Valparaíso\\
OII464: Desarrollo de Diseños Híbridos para Optimización\\
Profesor Emanuel Vega}
\end{center}

\begin{abstract}
\noindent
La planificación de cirugías electivas exige coordinar quirófanos y personal especializado,
respetando la precedencia entre anestesia y cirugía. Las metaheurísticas aplicadas a este
problema suelen fijar sus parámetros antes de la ejecución. Por su parte, el control aprendido
de parámetros se ha estudiado sobre todo para elegir operadores en algoritmos evolutivos.
Proponemos una arquitectura híbrida que añade a la optimización discreta por enjambre de
partículas (dPSO) un control aprendido de la perturbación en dos niveles. En el primero, un
bandido multibrazo elige cada cinco iteraciones la probabilidad de mutar una partícula. Este
nivel se actualiza con la mejora de aptitud del bloque. En el segundo, otro bandido elige en
cada mutación el movimiento aplicado: intercambio, inserción o reasignación de sala. Este
nivel se actualiza con el resultado de la partícula mutada. La configuración experimental
separa 37 instancias sintéticas en verificación, desarrollo y prueba. Comparamos siete
variantes con el mismo evaluador y $10\,000$ evaluaciones por corrida, incluida una ablación
de los dos niveles. El objetivo es establecer si aprender cuánto y cómo perturbar mejora la
calidad o la estabilidad de las soluciones, y a qué costo.
\end{abstract}

\section{Introducción}

Programar una jornada de cirugías electivas exige decidir el orden de atención y los quirófanos
asignados a cada procedimiento. Para comenzar una operación deben coincidir la disponibilidad
de sala y de personal elegible y la finalización de la etapa anterior. Cuando la cirugía se retrasa
después de la anestesia, aparecen esperas y retención de recursos que afectan el calendario
completo. Por ello, una planificación compacta debe considerar la interacción entre secuencias,
salas y personal. También debe considerar los tiempos auxiliares de preparación, transición y
limpieza.

Trabajamos con el repositorio docente del curso. Cada trabajo tiene dos operaciones, anestesia
y cirugía, con duraciones y recursos definidos en una instancia sintética. La decisión consiste en
una secuencia de prioridad y una asignación de sala por operación. El evaluador determina el
personal y los tiempos de atención. Cada instancia describe una jornada electiva determinista,
con todos los trabajos conocidos al inicio.

En este problema, cambiar el orden de atención altera la competencia por profesionales.
Reasignar una sala, en cambio, modifica la distribución de carga. La frecuencia y el tipo de esas
perturbaciones influyen en la búsqueda. Aplicarlas poco limita la exploración; aplicarlas
demasiado interrumpe el progreso. Además, un mismo movimiento no tiene por qué ser igual de
útil al comienzo y al final de una corrida. Sin embargo, las metaheurísticas aplicadas a la
programación quirúrgica fijan estos parámetros antes de ejecutar. Por otra parte, el control
aprendido se ha estudiado para elegir operadores en algoritmos evolutivos. No se ha usado para
regular cuánto se perturba un enjambre discreto (Sección~\ref{sec:arte}).

Nuestro objetivo es diseñar y evaluar un control de la perturbación que se añade a dPSO. El
control aprende, durante la ejecución, la intensidad y el tipo de esa perturbación. Las
contribuciones son tres:
\begin{enumerate}
\setlength{\itemsep}{0pt}
\item Un control en dos niveles, que aprende la probabilidad de mutación por bloques de
iteraciones y el movimiento en cada aplicación.
\item Una asignación de recompensa propia de cada nivel: agregada para la intensidad e
individual para el movimiento.
\item Un diseño experimental con ablación, que separa el efecto de añadir perturbación del
efecto de aprender a elegirla.
\end{enumerate}

\section{Estado del arte}
\label{sec:arte}

\paragraph{Programación de quirófanos.}
Cardoen et~al.~\cite{cardoen2010} y Samudra et~al.~\cite{samudra2016} ordenan el área según el
tipo de decisión, los recursos y las medidas de desempeño. En el nivel operativo, que es el de
este trabajo, los enfoques difieren en cómo tratan los recursos compartidos. Pham y
Klinkert~\cite{pham2008} formulan el caso quirúrgico como un taller de trabajos multimodo con
bloqueo. Plantean un programa entero mixto y lo prueban con un resolutor comercial. Bai
et~al.~\cite{bai2022} también parten de un modelo de taller, con tres estaciones, pero lo
resuelven con un algoritmo genético. Baesler et~al.~\cite{baesler2015}, en cambio, evalúan los
calendarios con un modelo de simulación. Un recocido simulado conectado a ese modelo busca
mejores secuencias de pacientes. A esta última línea pertenece el repositorio del curso. Su ventaja es que admite reglas de
asignación difíciles de linealizar. Su costo es que la calidad del resultado depende de la
metaheurística que propone los candidatos. En estos trabajos, los parámetros de la búsqueda se
fijan antes de ejecutar. Ninguno los revisa con la información que produce la corrida.

\paragraph{Control de parámetros.}
Eiben et~al.~\cite{eiben1999} revisan la terminología del área y clasifican los mecanismos que
modifican los parámetros durante la ejecución. Las revisiones posteriores de Karafotias
et~al.~\cite{karafotias2015} y de Aleti y Moser~\cite{aleti2016} tratan ese control en
algoritmos evolutivos. La segunda propone un modelo que divide la adaptación de un parámetro en
cuatro pasos. Los trabajos sobre selección de operadores difieren en cómo detectan y siguen al
mejor operador. Thierens~\cite{thierens2005} propone una estrategia que reacciona con rapidez a
los cambios y selecciona con alta probabilidad el mejor operador del momento. DaCosta
et~al.~\cite{dacosta2008} tratan la elección como un bandido multibrazo. Lo combinan con la
prueba estadística de Page-Hinkley para detectar cambios. Fialho et~al.~\cite{fialho2010}
analizan y comparan estos mecanismos. Subrayan que el problema es dinámico, mientras que los
bandidos clásicos suponen un entorno estático. Li et~al.~\cite{li2014} registran en una ventana
deslizante las tasas de mejora recientes de cada operador y basan la recompensa en su rango.
Todos estos trabajos eligen \emph{qué} operador aplicar. Ninguno aprende la \emph{intensidad}
con que se perturba.

En la optimización por enjambre de partículas, el control ha seguido otra vía. Shi y
Eberhart~\cite{shi2001} ajustan la inercia con reglas difusas. Ratnaweera
et~al.~\cite{ratnaweera2004} programan los coeficientes de aceleración y el tamaño de una
mutación en función del tiempo. Zhan et~al.~\cite{zhan2009} clasifican el estado del enjambre
y aplican reglas distintas en cada estado. Estos controles siguen un programa o unas
reglas fijadas de antemano. Los dos últimos se evaluaron con funciones de prueba. El aprendizaje por refuerzo dentro
de estos enjambres es más reciente~\cite{samma2016}. Las versiones discretas derivan de la
propuesta binaria de Kennedy y Eberhart~\cite{kennedy1997}. Heredan esos parámetros, pero no
un mecanismo de control equivalente. Por último, Talbi~\cite{talbi2021} y Karimi-Mamaghan et~al.~\cite{karimi2022} clasifican las
formas de integrar aprendizaje automático en una metaheurística. La segunda revisión incluye
entre ellas el ajuste de parámetros.

\paragraph{Brecha.}
De este contraste se desprenden tres vacíos. Primero, el control mediante bandidos se ha
aplicado a la elección entre operadores, no a la intensidad de la perturbación. Segundo, el
control revisado en enjambres de partículas sigue programas o reglas y se ha validado con
funciones de prueba. No se sabe, por tanto, si la
retroalimentación de la búsqueda es aprovechable en una versión discreta con secuencias y
asignaciones. Tercero, en la bibliografía revisada no encontramos un estudio que evalúe estos
mecanismos en programación quirúrgica con un evaluador de calendarios. Tampoco encontramos uno
que separe, con el mismo presupuesto, el efecto de añadir perturbación del efecto de aprender a
elegirla. Nuestra propuesta aborda los tres puntos. Aprende la intensidad y el movimiento con
señales de recompensa distintas, y se evalúa con una ablación que aísla cada nivel.

\section{Propuesta}

\subsection{Problema, representación y función de evaluación}

Sean $J$ el conjunto de trabajos, $R$ el de salas y $K=\{1,2\}$ el de operaciones. Un candidato
$x=(\pi,\rho)$ contiene una permutación $\pi$ de $J$ y una asignación $\rho_{jk}\in R_{jk}$. El
conjunto $R_{jk}\subseteq R$ reúne las salas elegibles para la operación $k$ del trabajo $j$.
Las dos operaciones de un trabajo pueden utilizar salas diferentes. El evaluador recorre los
trabajos según $\pi$. A cada operación le asigna el profesional elegible que se libera antes.

Minimizamos la función compuesta del evaluador del repositorio:
\begin{equation}
f(x)=\Cmax+\alpha\sum_{j\in J}\sum_{k\in K}S_{jk}+\beta\sum_{j\in J}W_j
+\gamma\,\maxop_{j\in J}W_j+\delta\,\Psi .
\label{eq:f}
\end{equation}
Sus términos son los siguientes. $\Cmax$ es el término de la última cirugía, incluida su
limpieza. $S_{jk}$ es el inicio de preparación de cada operación. $W_j$ es el intervalo entre
el fin de la anestesia y el inicio efectivo de la cirugía. $\Psi$ suma el coeficiente de
variación del número de operaciones por sala y la penalización $m(m+1)$ por $m$ salas sin uso.
Conservamos los pesos del repositorio: $\alpha=10^{-6}$, $\beta=0\dc5$, $\gamma=1\dc4$ y
$\delta=100$. Además de $f$, informaremos por separado $\Cmax$ y las esperas. La razón es que
una reducción de $f$ puede deberse solo al balance de salas.

\subsection{Componentes y flujo de la arquitectura híbrida}

La arquitectura combina dos componentes (Figura~\ref{fig:arq}). El componente metaheurístico
es dPSO. Actualiza las partículas a partir de sus velocidades y de las mejores soluciones
personal y global. El componente de aprendizaje consta de dos bandidos multibrazo no
contextuales (\emph{multi-armed bandits}), un método de aprendizaje por
recompensas~\cite{sutton2018}. El primero controla la intensidad: sus tres acciones son
probabilidades $p$ de mutar una partícula. El segundo controla el movimiento y también tiene
tres acciones. La primera intercambia dos posiciones de $\pi$. La segunda extrae un trabajo de
$\pi$ y lo inserta en otra posición. La tercera cambia la sala de una operación por otra
elegible. A diferencia de los trabajos de la Sección~\ref{sec:arte}, la arquitectura aprende a
la vez cuánto y cómo perturbar, y lo hace sobre un enjambre discreto.

\begin{figure}[t]
\centering
\begin{tikzpicture}[font=\small,>=Stealth,
  box/.style={draw,align=center,minimum height=1.25cm,inner sep=5pt}]
\node[box] (int) {Nivel de intensidad\\aprende $Q(p)$ y elige $p$\\\footnotesize cada $L$ iteraciones};
\node[box,below=0.55cm of int] (mov) {Nivel de movimiento\\aprende $V(m)$ y elige $m$\\\footnotesize en cada mutación};
\node[box,right=2.2cm of int,yshift=-0.9cm,minimum height=1.6cm] (pso) {dPSO + mutación\\partículas $(\pi,\rho)$};
\node[box,right=2.0cm of pso,minimum height=1.6cm] (ev) {Evaluador común\\aptitud y calendario};
\draw[->] (int.east) -- node[above,sloped]{$p$} (pso.150);
\draw[->] (mov.east) -- node[below,sloped]{$m$} (pso.210);
\draw[->] ([yshift=4pt]pso.east) -- node[above]{candidato $x$} ([yshift=4pt]ev.west);
\draw[->] ([yshift=-4pt]ev.west) -- node[below]{aptitud} ([yshift=-4pt]pso.east);
\draw[->] (pso.north) |- node[pos=0.75,above]{mejora del bloque} ([yshift=5mm]int.north) -- (int.north);
\draw[->] (pso.south) |- node[pos=0.75,below]{mejora de la partícula} ([yshift=-5mm]mov.south) -- (mov.south);
\end{tikzpicture}
\caption{Intercambio entre los dos niveles de control, dPSO y el evaluador. La probabilidad $p$
se mantiene durante un bloque y se recompensa con la mejora de la mejor solución; el
movimiento $m$ se elige en cada mutación y se recompensa con la mejora de la partícula mutada.}
\label{fig:arq}
\end{figure}

La mutación actúa después de la actualización habitual de dPSO y antes de evaluar la nueva
posición. En ese punto, cada partícula se muta con probabilidad $p$ mediante el movimiento $m$
elegido. Los tres movimientos preservan la permutación y la elegibilidad de salas; el evaluador
comprueba la factibilidad temporal. Si un movimiento no está disponible, elegimos al azar entre
los restantes. Anulamos las velocidades en las coordenadas modificadas, con el mismo criterio en
todas las variantes. La posición actual puede empeorar. Las memorias personal y global conservan
las mejores soluciones encontradas.

\subsection{Nivel de intensidad}

Las acciones son $\mathcal{P}=\{0\dc05;\,0\dc15;\,0\dc30\}$. Al comienzo de cada bloque de $L=5$
iteraciones, el control elige una acción $p_t$, que permanece fija durante el bloque. Los
valores $Q_t(p)$ se inicializan en cero y las tres acciones se prueban una vez en orden
aleatorio. Después aplicamos una política épsilon-voraz. Con probabilidad $1-\epsilon=0\dc9$
elige una acción de mayor $Q_t$; con probabilidad $\epsilon=0\dc1$, una acción uniforme. Los
empates se resuelven al azar. Sean $g_t$ y $g_{t+1}$ las mejores aptitudes antes y después del
bloque. La recompensa y la actualización son
\begin{align}
r_t &= \minop\Bigl\{1,\;\maxop\Bigl\{0,\;\frac{g_t-g_{t+1}}{\maxop\{1,|g_t|\}}\Bigr\}\Bigr\},
\label{eq:r}\\
Q_{t+1}(p_t) &= Q_t(p_t)+\eta\bigl(r_t-Q_t(p_t)\bigr),\qquad \eta=0\dc2 .
\label{eq:q}
\end{align}
Por ejemplo, si la mejor aptitud baja de 1000 a 980, entonces $r_t=0\dc02$. En ese caso, un
valor $Q_t(p_t)=0\dc010$ pasa a $0\dc012$. El paso constante da más peso a las observaciones
recientes~\cite{sutton2018}. Esto es adecuado porque la utilidad de cada intensidad cambia a lo
largo de la búsqueda~\cite{fialho2010}.

\subsection{Nivel de movimiento}

Las acciones son $\mathcal{M}=\{\text{intercambio},\text{inserción},\text{sala}\}$. Cada vez que
una partícula $i$ va a ser mutada, el control elige $m$ con la misma política épsilon-voraz
($\epsilon=0\dc1$). La política usa los valores $V(m)$, inicializados en cero. Tras evaluar la
posición mutada, la recompensa es binaria. La actualización usa un paso menor, porque este
nivel recibe muchas más observaciones:
\begin{equation}
u_i=\begin{cases}1 & \text{si } f(x_i^{\text{nueva}})<f(x_i^{\text{anterior}}),\\ 0 & \text{en otro caso,}\end{cases}
\qquad V(m)\leftarrow V(m)+\eta_m\bigl(u_i-V(m)\bigr),\quad \eta_m=0\dc05 .
\label{eq:v}
\end{equation}
Así, $V(m)$ estima la frecuencia reciente con que una partícula mejora cuando se le aplica $m$.
Los dos niveles usan señales distintas a propósito. La intensidad afecta a todo el enjambre
durante un bloque. Por eso la recompensamos con la mejora de la mejor solución (unas 66
decisiones por corrida). El movimiento afecta a una sola partícula. Por eso lo recompensamos
con el resultado inmediato de esa partícula (del orden de 1500 observaciones por corrida). En
ambos casos, la señal incluye el efecto del paso de dPSO que precede a la mutación. Las
variantes de comparación de la Sección~\ref{sec:exp} miden si, pese a ese ruido, la señal es
útil.

\subsection{Integración, ejecución y recursos}

Antes de ejecutar solo validamos las instancias y fijamos los parámetros del control. Todo el
aprendizaje ocurre durante la optimización de cada instancia. Los valores $Q$ y $V$ se
reinician en cada corrida y no se transfiere información entre instancias. Si la población
inicial no contiene ninguna solución factible, usamos $p=0\dc15$ hasta obtenerla. El criterio
de parada es agotar $B=10\,000$ llamadas al evaluador, incluida la evaluación inicial. El último
bloque se trunca al presupuesto restante y no actualiza $Q$. El control almacena seis
estimaciones. Realiza una actualización escalar por bloque y otra por mutación, con memoria
constante.

\section{Configuración experimental}
\label{sec:exp}

\paragraph{Preguntas.}
El experimento responde tres preguntas. P1: ?`mejora dPSO al añadir la mutación, con
independencia de cómo se controle? P2: ?`aporta el aprendizaje de la intensidad algo más que
una probabilidad fija bien elegida, una elección aleatoria o un ajuste por reglas? P3: ?`aporta
el aprendizaje del movimiento, solo o combinado con el de la intensidad?

\paragraph{Instancias.}
Usamos las 37 instancias sintéticas del repositorio del curso, de 3 a 60 trabajos y de 2 a 12
salas. Cada una se generó con una semilla fija y con distribuciones de duración documentadas en
el propio archivo. No contienen datos de pacientes. Todas las salas son elegibles para cada
operación; el personal elegible es un subconjunto por operación. Usamos las instancias sin
preprocesamiento. La Tabla~\ref{tab:inst} muestra su división. Las instancias de verificación
solo sirven para comprobar la implementación. Las de desarrollo sirven para fijar los
parámetros. Las de prueba no se ejecutan hasta que la configuración está cerrada. Nueve
instancias de prueba (40 a 60 trabajos) quedan fuera del rango de tamaños de desarrollo.
Permiten observar el comportamiento al aumentar la escala.

\begin{table}[htb]
\centering\small
\caption{División de las instancias. STD: familia estándar (6 a 9 salas); 12S: familia de doce salas.}
\label{tab:inst}
\begin{tabular}{@{}llcl@{}}
\toprule
Conjunto & Instancias & $n$ & Uso \\
\midrule
Verificación & didácticas de 3, 5, 8 y 10 trabajos & 4 & restricciones y contador de evaluaciones \\
Desarrollo & réplica 01 de STD y 12S, 15 a 30 trabajos & 8 & elección de $p$ fija, piloto de tiempos \\
Prueba & réplicas 02 y 03 de STD y 12S, 15 a 30 trabajos & 16 & comparación final \\
       & réplicas 01 a 03 de 12S, 40 a 60 trabajos & 9 & comparación final, mayor escala \\
\bottomrule
\end{tabular}
\end{table}

\paragraph{Métodos comparados.}
Comparamos siete variantes de dPSO (Tabla~\ref{tab:met}). Para que la comparación sea justa,
todas comparten el evaluador, el presupuesto y la población inicial de cada semilla. Las que
mutan comparten además los tres movimientos y el tratamiento de velocidades y memorias. A1 es
la dPSO del repositorio. A2 usa el valor de $\mathcal{P}$ con menor rango promedio en
desarrollo. A3 comparte el arranque de A5 y después elige $p$ uniformemente, sin usar $Q$. A4
empieza en $0\dc15$ y revisa $p$ al terminar cada bloque. Baja un nivel si la mejor aptitud
mejoró más de $10^{-6}$ en términos relativos; si no, sube uno. A5 es la propuesta del primer
informe, A6 aprende solo el movimiento y A7 es la propuesta completa. A2, A5, A6 y A7 forman un
diseño $2\times2$: intensidad fija o aprendida, movimiento uniforme o aprendido. Cada pregunta
tiene sus contrastes. P1 se responde con A1 frente a A2. P2, con A5 frente a A2, A3 y A4. P3,
con A6 frente a A2 y con A7 frente a A5.

\begin{table}[htb]
\centering\small
\caption{Variantes comparadas. Todas usan 30 partículas y $B=10\,000$ evaluaciones.}
\label{tab:met}
\begin{tabular}{@{}llll@{}}
\toprule
Variante & Intensidad $p$ & Movimiento $m$ & Papel \\
\midrule
A1 & sin mutación & --- & línea base \\
A2 & fija ($0\dc30$) & uniforme & mutación sin control \\
A3 & aleatoria por bloque & uniforme & control sin información \\
A4 & reglas & uniforme & control adaptativo sin aprendizaje \\
A5 & aprendida & uniforme & ablación: solo intensidad \\
A6 & fija ($0\dc30$) & aprendido & ablación: solo movimiento \\
A7 & aprendida & aprendido & propuesta completa \\
\bottomrule
\end{tabular}
\end{table}

\paragraph{Parámetros.}
Los parámetros de dPSO son los del repositorio: 30 partículas, inercia $w=0\dc7$, coeficientes
$c_1=c_2=1\dc5$ y velocidades en $[-4,4]$. Los del control son
$\mathcal{P}=\{0\dc05;\,0\dc15;\,0\dc30\}$, $L=5$, $\epsilon=0\dc1$, $\eta=0\dc2$ y
$\eta_m=0\dc05$. Fijamos estos valores a priori, con los habituales para pasos constantes y
exploración épsilon-voraz~\cite{sutton2018}. No los ajustamos. Con solo ocho instancias de
desarrollo, una calibración fina favorecería a las variantes aprendidas frente a las de
referencia. El único valor elegido con datos es la probabilidad de A2 y A6. Para ello
ejecutamos un piloto con las ocho instancias de desarrollo y 20 semillas. Los rangos promedio
de la aptitud mediana fueron $1\dc50$ para $p=0\dc30$, $2\dc13$ para $p=0\dc15$ y $2\dc38$ para
$p=0\dc05$. Por tanto, fijamos $p=0\dc30$. Como análisis complementario, mediremos en desarrollo
la sensibilidad de A7 a $L\in\{3,5,10\}$ y $\epsilon\in\{0\dc05;\,0\dc1;\,0\dc2\}$.

\paragraph{Protocolo de ejecución.}
Ejecutamos cada variante 30 veces por instancia, con las semillas 1 a 30, comunes a todas las
variantes. La semilla determina cuatro flujos aleatorios independientes: población inicial,
dPSO, mutación y control. Así, dos variantes con la misma semilla parten del mismo enjambre.
Toda llamada al evaluador cuenta en el presupuesto, y cada partícula se evalúa una vez por
iteración. En consecuencia, una corrida comprende la evaluación inicial, 66 bloques completos y
un bloque final truncado de 70 evaluaciones. La prueba final suma $25\times7\times30=5250$
corridas y $52\dc5$ millones de evaluaciones.

\paragraph{Verificación previa.}
Verificamos la implementación con las cuatro instancias didácticas. Cada corrida consume
exactamente $10\,000$ evaluaciones y el resultado es reproducible con la misma semilla. Los
calendarios devueltos respetan la precedencia entre anestesia y cirugía. Tampoco solapan
operaciones en una misma sala ni en un mismo profesional. En las 1440 corridas del piloto, todas
las soluciones fueron factibles. La causa está en el evaluador: aplica como límite de espera el
mayor entre el de la instancia y una cota global de $10\,000$ minutos. Por ello, el límite
propio de cada instancia (27 a 53 minutos) no actúa como restricción. Las esperas solo se
penalizan mediante $\beta$ y $\gamma$. Mantenemos esta regla en todas las variantes e
informaremos la espera máxima para que el efecto sea visible.

\paragraph{Entorno.}
La implementación está en Python~3.13 con NumPy~2.4. Usa la entrada que el repositorio reserva
para la metaheurística de los alumnos, por lo que conserva su interfaz y su evaluador. Las
corridas se ejecutan en CPU, en un equipo con procesador AMD Ryzen~5~9600 (6 núcleos, 12 hilos).
El equipo tiene 32~GB de memoria y Windows~11. Usamos 12 procesos en paralelo, con una corrida
por proceso. En el piloto, y con el equipo cargado, una corrida tardó entre $2\dc4$~s (15
trabajos) y $4\dc0$~s (30 trabajos). No hubo diferencias apreciables entre variantes. Estimamos
menos de una hora para la prueba final. Junto con los resultados registramos las semillas, los
parámetros, la versión del intérprete y el equipo.

\paragraph{Medidas.}
Por corrida registramos la aptitud final $f$, $\Cmax$, la espera total y la máxima, el tiempo de
ejecución y el número de mutaciones. Para las variantes con control guardamos, por bloque, la
acción elegida, la recompensa y los valores $Q$ y $V$. También guardamos, por corrida, el número
de usos de cada movimiento. Resumimos la calidad por instancia con la mediana y el rango
intercuartílico de las 30 corridas. En el piloto, ese rango fue en promedio el 17~\% de la
mediana. Esta dispersión justifica usar medianas y 30 repeticiones. Las curvas de mejor aptitud
frente al número de evaluaciones describen el progreso. La frecuencia de cada acción a lo largo
de la corrida describe lo que aprende el control.

\paragraph{Análisis estadístico.}
La unidad de comparación es la instancia. Para cada variante e instancia calculamos la mediana
de $f$ en las 30 corridas. Sobre las 25 instancias de prueba aplicamos la prueba de Friedman a
los rangos de esas medianas. Si se rechaza la igualdad, comparamos A7 con cada una de las otras
seis variantes. Usamos la prueba de rangos con signo de Wilcoxon y la corrección de Holm, con
un nivel de significación del 5~\%~\cite{derrac2011}. Algunos contrastes de P1 a P3 no incluyen
a A7: A1 con A2; A5 con A2, A3 y A4; y A6 con A2. Los analizamos del mismo modo, como una
segunda familia. Junto con cada valor $p$ informamos la diferencia relativa mediana. También
informamos el número de instancias ganadas, empatadas y perdidas. Desglosamos los resultados
por tamaño (15 a 30 y 40 a 60 trabajos) y repetimos el análisis para $\Cmax$.

\section{Conclusión y trabajo futuro}

El problema abordado exige coordinar el orden de atención y las salas bajo recursos
compartidos. La literatura ha aprendido qué operador aplicar en algoritmos evolutivos. También
ha ajustado con reglas los coeficientes de los enjambres de partículas. No ha estudiado, en cambio,
el aprendizaje conjunto de la intensidad y del tipo de perturbación en un enjambre discreto
aplicado a la programación quirúrgica. Para cubrir esa brecha proponemos un control en dos
niveles, con una recompensa adaptada a cada uno. La configuración experimental separa
instancias de desarrollo y de prueba, iguala el presupuesto de evaluaciones y aísla cada nivel
mediante ablación. La implementación y un piloto en desarrollo confirman que el protocolo es
ejecutable en menos de una hora de cómputo. Los resultados de la comparación se presentarán en
la siguiente entrega.

\paragraph{Limitaciones.}
El trabajo tiene tres limitaciones. Primero, las instancias son sintéticas y deterministas.
Las conclusiones no se extenderán, por tanto, a datos hospitalarios ni a duraciones inciertas.
Segundo, la comparación es interna a dPSO: mide el aporte del control, no la posición de dPSO
frente a otras metaheurísticas. Tercero, la función objetivo agrega cinco términos con pesos
fijos, y el balance de salas pesa mucho en ella. Por eso informamos también $\Cmax$ y las
esperas.

\paragraph{Trabajo futuro.}
El paso inmediato es ejecutar la prueba final y analizar si el aprendizaje mejora la calidad o
la estabilidad de las soluciones. Después vemos tres extensiones. La primera es dar contexto al
control, por ejemplo con una medida del estancamiento del enjambre. La segunda es comparar con
las otras metaheurísticas del repositorio bajo el mismo presupuesto. La tercera es evaluar
instancias con duraciones inciertas. En una perspectiva más amplia, un control que aprende
durante la ejecución podría evitar la calibración manual para cada hospital. Esa propiedad
acercaría estas técnicas a la planificación quirúrgica real.

\begin{thebibliography}{99}\small
\setlength{\itemsep}{2pt}

\bibitem{cardoen2010} Cardoen, B.; Demeulemeester, E.; Beliën, J. Operating room planning and
scheduling: A literature review. \emph{European Journal of Operational Research} \textbf{2010},
201(3), 921--932. doi:10.1016/j.ejor.2009.04.011.

\bibitem{samudra2016} Samudra, M.; Van Riet, C.; Demeulemeester, E.; Cardoen, B.; Vansteenkiste, N.;
Rademakers, F.E. Scheduling operating rooms: achievements, challenges and pitfalls.
\emph{Journal of Scheduling} \textbf{2016}, 19(5), 493--525. doi:10.1007/s10951-016-0489-6.

\bibitem{pham2008} Pham, D.-N.; Klinkert, A. Surgical case scheduling as a generalized job shop
scheduling problem. \emph{European Journal of Operational Research} \textbf{2008}, 185(3),
1011--1025. doi:10.1016/j.ejor.2006.03.059.

\bibitem{bai2022} Bai, X.; Zhang, W.; Luo, L.; Ma, H.; Zhu, T. Day Surgery Scheduling and
Optimization in Large Public Hospitals in China: A Three-Station Job Shop Scheduling Problem.
\emph{Journal of Healthcare Engineering} \textbf{2022}, 2022, 1149657. doi:10.1155/2022/1149657.

\bibitem{baesler2015} Baesler, F.; Gatica, J.; Correa, R. Simulation optimisation for operating room
scheduling. \emph{International Journal of Simulation Modelling} \textbf{2015}, 14(2), 215--226.
doi:10.2507/IJSIMM14(2)3.287.

\bibitem{eiben1999} Eiben, A.E.; Hinterding, R.; Michalewicz, Z. Parameter control in evolutionary
algorithms. \emph{IEEE Transactions on Evolutionary Computation} \textbf{1999}, 3(2), 124--141.
doi:10.1109/4235.771166.

\bibitem{karafotias2015} Karafotias, G.; Hoogendoorn, M.; Eiben, A.E. Parameter Control in
Evolutionary Algorithms: Trends and Challenges. \emph{IEEE Transactions on Evolutionary
Computation} \textbf{2015}, 19(2), 167--187. doi:10.1109/TEVC.2014.2308294.

\bibitem{aleti2016} Aleti, A.; Moser, I. A Systematic Literature Review of Adaptive Parameter
Control Methods for Evolutionary Algorithms. \emph{ACM Computing Surveys} \textbf{2016}, 49(3),
1--35. doi:10.1145/2996355.

\bibitem{thierens2005} Thierens, D. An adaptive pursuit strategy for allocating operator
probabilities. En \emph{Proceedings of the 7th Annual Conference on Genetic and Evolutionary
Computation}; ACM, \textbf{2005}; pp.~1539--1546. doi:10.1145/1068009.1068251.

\bibitem{dacosta2008} DaCosta, L.; Fialho, Á.; Schoenauer, M.; Sebag, M. Adaptive operator
selection with dynamic multi-armed bandits. En \emph{Proceedings of the 10th Annual Conference on
Genetic and Evolutionary Computation}; ACM, \textbf{2008}; pp.~913--920.
doi:10.1145/1389095.1389272.

\bibitem{fialho2010} Fialho, Á.; Da Costa, L.; Schoenauer, M.; Sebag, M. Analyzing bandit-based
adaptive operator selection mechanisms. \emph{Annals of Mathematics and Artificial Intelligence}
\textbf{2010}, 60(1--2), 25--64. doi:10.1007/s10472-010-9213-y.

\bibitem{li2014} Li, K.; Fialho, Á.; Kwong, S.; Zhang, Q. Adaptive Operator Selection With Bandits
for a Multiobjective Evolutionary Algorithm Based on Decomposition. \emph{IEEE Transactions on
Evolutionary Computation} \textbf{2014}, 18(1), 114--130. doi:10.1109/TEVC.2013.2239648.

\bibitem{shi2001} Shi, Y.; Eberhart, R.C. Fuzzy adaptive particle swarm optimization. En
\emph{Proceedings of the 2001 Congress on Evolutionary Computation}; IEEE, \textbf{2001}; vol.~1,
pp.~101--106. doi:10.1109/CEC.2001.934377.

\bibitem{ratnaweera2004} Ratnaweera, A.; Halgamuge, S.K.; Watson, H.C. Self-Organizing Hierarchical
Particle Swarm Optimizer With Time-Varying Acceleration Coefficients. \emph{IEEE Transactions on
Evolutionary Computation} \textbf{2004}, 8(3), 240--255. doi:10.1109/TEVC.2004.826071.

\bibitem{zhan2009} Zhan, Z.-H.; Zhang, J.; Li, Y.; Chung, H.S.-H. Adaptive Particle Swarm
Optimization. \emph{IEEE Transactions on Systems, Man, and Cybernetics, Part B (Cybernetics)}
\textbf{2009}, 39(6), 1362--1381. doi:10.1109/TSMCB.2009.2015956.

\bibitem{samma2016} Samma, H.; Lim, C.P.; Mohamad Saleh, J. A new Reinforcement Learning-based
Memetic Particle Swarm Optimizer. \emph{Applied Soft Computing} \textbf{2016}, 43, 276--297.
doi:10.1016/j.asoc.2016.01.006.

\bibitem{kennedy1997} Kennedy, J.; Eberhart, R.C. A discrete binary version of the particle swarm
algorithm. En \emph{1997 IEEE International Conference on Systems, Man, and Cybernetics}; IEEE,
\textbf{1997}; vol.~5, pp.~4104--4108. doi:10.1109/ICSMC.1997.637339.

\bibitem{talbi2021} Talbi, E.-G. Machine Learning into Metaheuristics: A Survey and Taxonomy.
\emph{ACM Computing Surveys} \textbf{2021}, 54(6), 1--32. doi:10.1145/3459664.

\bibitem{karimi2022} Karimi-Mamaghan, M.; Mohammadi, M.; Meyer, P.; Karimi-Mamaghan, A.M.;
Talbi, E.-G. Machine learning at the service of meta-heuristics for solving combinatorial
optimization problems: A state-of-the-art. \emph{European Journal of Operational Research}
\textbf{2022}, 296(2), 393--422. doi:10.1016/j.ejor.2021.04.032.

\bibitem{sutton2018} Sutton, R.S.; Barto, A.G. \emph{Reinforcement Learning: An Introduction},
2.\textsuperscript{a}~ed.; MIT Press, \textbf{2018}; cap.~2. ISBN: 9780262039246.

\bibitem{derrac2011} Derrac, J.; García, S.; Molina, D.; Herrera, F. A practical tutorial on the use
of nonparametric statistical tests as a methodology for comparing evolutionary and swarm
intelligence algorithms. \emph{Swarm and Evolutionary Computation} \textbf{2011}, 1(1), 3--18.
doi:10.1016/j.swevo.2011.02.002.

\end{thebibliography}

\end{document}
